% FILE: CSE-24_TermFinal.tex
\documentclass{article}
\usepackage{amsmath}
\usepackage{amssymb}
\begin{document}

%%%%%%%%%%%%%%%%%%%%%%%%%%%%%%%%%%%%%%%%%%%%%%%%%%
% QUESTION_START
%%%%%%%%%%%%%%%%%%%%%%%%%%%%%%%%%%%%%%%%%%%%%%%%%%
% id=cse24-final-seca-q1a
% discipline=CSE
% year=2024
% exam_type=Term Final
% section=A
% ct_number=UNKNOWN
% question_number=1a
% topic=Continuity
% subtopic=General
% difficulty=Medium
% length=Long
% question_type=New
% frequency=1
% appearances=
% OCR_STATUS=VERIFIED
\section*{Term Final (Sec A) - Q1(a)}
Question:
Define function, even function and odd function with examples. Draw the graph and also find the domain and range of $f(x) = |x-4| + |x| + |x+6|$.

Solution:
Step 1: Define Function, Even Function, and Odd Function.
\begin{itemize}
    \item \textbf{Function}: A relation $f$ from a set $A$ (domain) to a set $B$ (codomain) is a function if every element $x \in A$ is associated with exactly one element $y = f(x) \in B$.
    \\ \textit{Example}: $f(x) = x^2$.
    \item \textbf{Even Function}: A function $f(x)$ is even if $f(-x) = f(x)$ for all $x$ in its domain. Its graph is symmetric with respect to the $y$-axis.
    \\ \textit{Example}: $f(x) = \cos x$.
    \item \textbf{Odd Function}: A function $f(x)$ is odd if $f(-x) = -f(x)$ for all $x$ in its domain. Its graph is symmetric with respect to the origin.
    \\ \textit{Example}: $f(x) = \sin x$.
\end{itemize}

Step 2: Determine the Domain of $f(x) = |x-4| + |x| + |x+6|$.
Since $f(x)$ is a sum of absolute value functions, which are defined for all real numbers:
\[
\text{Domain} = (-\infty, \infty) \quad \text{or} \quad \mathbb{R}
\]

Step 3: Define $f(x)$ piecewise to find the Range and graph behavior.
The critical points of the absolute values are $x = -6$, $x = 0$, and $x = 4$.
\begin{itemize}
    \item \textbf{Case 1: $x < -6$}
    \[
    f(x) = -(x-4) - x - (x+6) = -x+4-x-x-6 = -3x - 2
    \]
    \item \textbf{Case 2: $-6 \le x < 0$}
    \[
    f(x) = -(x-4) - x + (x+6) = -x+4-x+x+6 = -x + 10
    \]
    \item \textbf{Case 3: $0 \le x < 4$}
    \[
    f(x) = -(x-4) + x + (x+6) = -x+4+x+x+6 = x + 10
    \]
    \item \textbf{Case 4: $x \ge 4$}
    \[
    f(x) = (x-4) + x + (x+6) = 3x + 2
    \]
\end{itemize}

Step 4: Find the Range.
Evaluate the function at the critical points:
\begin{align*}
    f(-6) &= |-6-4| + |-6| + |-6+6| = 10 + 6 + 0 = 16 \\
    f(0) &= |0-4| + |0| + |0+6| = 4 + 0 + 6 = 10 \\
    f(4) &= |4-4| + |4| + |4+6| = 0 + 4 + 10 = 14
\end{align*}
The global minimum of the function occurs at $x = 0$ with value $10$. As $x \to \pm\infty$, $f(x) \to \infty$. Thus:
\[
\text{Range} = [10, \infty)
\]

Step 5: Describe the Graph.
The graph is a continuous, piecewise linear curve consisting of four segments:
\begin{itemize}
    \item A line with slope $-3$ for $x \in (-\infty, -6)$ ending at $(-6, 16)$.
    \item A line with slope $-1$ for $x \in [-6, 0)$ connecting $(-6, 16)$ to $(0, 10)$.
    \item A line with slope $+1$ for $x \in [0, 4)$ connecting $(0, 10)$ to $(4, 14)$.
    \item A line with slope $+3$ for $x \in [4, \infty)$ originating from $(4, 14)$.
\end{itemize}

Final Answer:
\[
\boxed{\text{Domain} = (-\infty, \infty), \quad \text{Range} = [10, \infty)}
\]
%%%%%%%%%%%%%%%%%%%%%%%%%%%%%%%%%%%%%%%%%%%%%%%%%%
% QUESTION_END
%%%%%%%%%%%%%%%%%%%%%%%%%%%%%%%%%%%%%%%%%%%%%%%%%%

%%%%%%%%%%%%%%%%%%%%%%%%%%%%%%%%%%%%%%%%%%%%%%%%%%
% QUESTION_START
%%%%%%%%%%%%%%%%%%%%%%%%%%%%%%%%%%%%%%%%%%%%%%%%%%
% id=cse24-final-seca-q1b
% discipline=CSE
% year=2024
% exam_type=Term Final
% section=A
% ct_number=UNKNOWN
% question_number=1b
% topic=Continuity
% subtopic=General
% difficulty=Medium
% length=Medium
% question_type=New
% frequency=1
% appearances=
% OCR_STATUS=VERIFIED
\section*{Term Final (Sec A) - Q1(b)}
Question:
Define limit of a function. A function is defined by
\[
f(x) = \begin{cases} 
x^2, & \text{when } x \le 1 \\ 
x, & \text{when } 1 < x \le 2 \\ 
\frac{1}{4}x^3, & \text{when } x > 2 
\end{cases}
\]
Find $\lim_{x \to 1} f(x)$ and discuss the continuity of $f(x)$ at $x = 1$.

Solution:
Step 1: Define the limit of a function.
Let $f(x)$ be defined on an open interval containing $c$, except possibly at $c$ itself. We say that the limit of $f(x)$ as $x$ approaches $c$ is $L$, written as:
\[
\lim_{x \to c} f(x) = L
\]
if for every $\epsilon > 0$ there exists a $\delta > 0$ such that $0 < |x - c| < \delta \implies |f(x) - L| < \epsilon$.

Step 2: Find the left-hand limit (LHL) at $x = 1$.
For $x \le 1$, $f(x) = x^2$:
\[
\text{LHL} = \lim_{x \to 1^-} f(x) = \lim_{x \to 1^-} x^2 = 1^2 = 1
\]

Step 3: Find the right-hand limit (RHL) at $x = 1$.
For $1 < x \le 2$, $f(x) = x$:
\[
\text{RHL} = \lim_{x \to 1^+} f(x) = \lim_{x \to 1^+} x = 1
\]

Step 4: Determine the existence of the limit.
Since $\text{LHL} = \text{RHL} = 1$:
\[
\lim_{x \to 1} f(x) = 1
\]

Step 5: Discuss continuity at $x = 1$.
A function is continuous at $x = c$ if:
\[
\lim_{x \to c} f(x) = f(c)
\]
Here, the functional value at $x = 1$ is:
\[
f(1) = 1^2 = 1
\]
Since $\lim_{x \to 1} f(x) = 1 = f(1)$, the function $f(x)$ is continuous at $x = 1$.

Final Answer:
\[
\boxed{\lim_{x \to 1} f(x) = 1, \quad \text{and the function is continuous at } x = 1.}
\]
%%%%%%%%%%%%%%%%%%%%%%%%%%%%%%%%%%%%%%%%%%%%%%%%%%
% QUESTION_END
%%%%%%%%%%%%%%%%%%%%%%%%%%%%%%%%%%%%%%%%%%%%%%%%%%

%%%%%%%%%%%%%%%%%%%%%%%%%%%%%%%%%%%%%%%%%%%%%%%%%%
% QUESTION_START
%%%%%%%%%%%%%%%%%%%%%%%%%%%%%%%%%%%%%%%%%%%%%%%%%%
% id=cse24-final-seca-q2ai
% discipline=CSE
% year=2024
% exam_type=Term Final
% section=A
% ct_number=UNKNOWN
% question_number=2ai
% topic=Differentiation
% subtopic=Implicit Differentiation
% difficulty=Medium
% length=Short
% question_type=New
% frequency=1
% appearances=
% OCR_STATUS=VERIFIED
\section*{Term Final (Sec A) - Q2(a)(i)}
Question:
Find $\frac{dy}{dx}$ of $(\cos x)^y = (\sin y)^x$.

Solution:
Step 1: Take the natural logarithm of both sides.
\[
\ln\left( (\cos x)^y \right) = \ln\left( (\sin y)^x \right) \implies y \ln(\cos x) = x \ln(\sin y)
\]

Step 2: Differentiate both sides implicitly with respect to $x$ using the product rule.
\[
\frac{d}{dx} [y \ln(\cos x)] = \frac{d}{dx} [x \ln(\sin y)]
\]
\[
\frac{dy}{dx} \ln(\cos x) + y \cdot \frac{1}{\cos x}(-\sin x) = 1 \cdot \ln(\sin y) + x \cdot \frac{1}{\sin y}(\cos y) \frac{dy}{dx}
\]
\[
\frac{dy}{dx} \ln(\cos x) - y \tan x = \ln(\sin y) + x \cot y \frac{dy}{dx}
\]

Step 3: Group the terms involving $\frac{dy}{dx}$.
\[
\frac{dy}{dx} \ln(\cos x) - x \cot y \frac{dy}{dx} = \ln(\sin y) + y \tan x
\]
\[
\frac{dy}{dx} \left[ \ln(\cos x) - x \cot y \right] = \ln(\sin y) + y \tan x
\]

Step 4: Solve for $\frac{dy}{dx}$.
\[
\frac{dy}{dx} = \frac{\ln(\sin y) + y \tan x}{\ln(\cos x) - x \cot y}
\]

Final Answer:
\[
\boxed{\frac{dy}{dx} = \frac{\ln(\sin y) + y \tan x}{\ln(\cos x) - x \cot y}}
\]
%%%%%%%%%%%%%%%%%%%%%%%%%%%%%%%%%%%%%%%%%%%%%%%%%%
% QUESTION_END
%%%%%%%%%%%%%%%%%%%%%%%%%%%%%%%%%%%%%%%%%%%%%%%%%%

%%%%%%%%%%%%%%%%%%%%%%%%%%%%%%%%%%%%%%%%%%%%%%%%%%
% QUESTION_START
%%%%%%%%%%%%%%%%%%%%%%%%%%%%%%%%%%%%%%%%%%%%%%%%%%
% id=cse24-final-seca-q2aii
% discipline=CSE
% year=2024
% exam_type=Term Final
% section=B
% ct_number=UNKNOWN
% question_number=2aii
% topic=Differentiation
% subtopic=Logarithmic Differentiation
% difficulty=Hard
% length=Medium
% question_type=New
% frequency=1
% appearances=
% OCR_STATUS=VERIFIED
\section*{Term Final (Sec A) - Q2(a)(ii)}
Question:
Find $\frac{dy}{dx}$ of $y = (\tan x)^{\cot x} + (\cot x)^{\tan x}$.

Solution:
Step 1: Express the terms separately.
Let $y = u + v$, where $u = (\tan x)^{\cot x}$ and $v = (\cot x)^{\tan x}$. Then $\frac{dy}{dx} = \frac{du}{dx} + \frac{dv}{dx}$.

Step 2: Find the derivative of $u = (\tan x)^{\cot x}$ using logarithmic differentiation.
\[
\ln u = \cot x \ln(\tan x)
\]
Differentiate with respect to $x$:
\[
\frac{1}{u} \frac{du}{dx} = -\csc^2 x \ln(\tan x) + \cot x \cdot \frac{1}{\tan x} \sec^2 x
\]
Simplify the second term:
\[
\cot x \cdot \cot x \cdot \sec^2 x = \frac{\cos^2 x}{\sin^2 x} \cdot \frac{1}{\cos^2 x} = \frac{1}{\sin^2 x} = \csc^2 x
\]
So:
\[
\frac{1}{u} \frac{du}{dx} = -\csc^2 x \ln(\tan x) + \csc^2 x = \csc^2 x [1 - \ln(\tan x)]
\]
\[
\frac{du}{dx} = (\tan x)^{\cot x} \csc^2 x [1 - \ln(\tan x)]
\]

Step 3: Find the derivative of $v = (\cot x)^{\tan x}$ using logarithmic differentiation.
\[
\ln v = \tan x \ln(\cot x)
\]
Differentiate with respect to $x$:
\[
\frac{1}{v} \frac{dv}{dx} = \sec^2 x \ln(\cot x) + \tan x \cdot \frac{1}{\cot x} (-\csc^2 x)
\]
Simplify the second term:
\[
\tan x \cdot \tan x \cdot (-\csc^2 x) = -\frac{\sin^2 x}{\cos^2 x} \cdot \frac{1}{\sin^2 x} = -\sec^2 x
\]
So:
\[
\frac{1}{v} \frac{dv}{dx} = \sec^2 x \ln(\cot x) - \sec^2 x = \sec^2 x [\ln(\cot x) - 1]
\]
\[
\frac{dv}{dx} = (\cot x)^{\tan x} \sec^2 x [\ln(\cot x) - 1]
\]

Step 4: Combine the derivatives.
\[
\frac{dy}{dx} = (\tan x)^{\cot x} \csc^2 x [1 - \ln(\tan x)] + (\cot x)^{\tan x} \sec^2 x [\ln(\cot x) - 1]
\]

Final Answer:
\[
\boxed{\frac{dy}{dx} = (\tan x)^{\cot x} \csc^2 x [1 - \ln(\tan x)] + (\cot x)^{\tan x} \sec^2 x [\ln(\cot x) - 1]}
\]
%%%%%%%%%%%%%%%%%%%%%%%%%%%%%%%%%%%%%%%%%%%%%%%%%%
% QUESTION_END
%%%%%%%%%%%%%%%%%%%%%%%%%%%%%%%%%%%%%%%%%%%%%%%%%%

%%%%%%%%%%%%%%%%%%%%%%%%%%%%%%%%%%%%%%%%%%%%%%%%%%
% QUESTION_START
%%%%%%%%%%%%%%%%%%%%%%%%%%%%%%%%%%%%%%%%%%%%%%%%%%
% id=cse24-final-seca-q2aiii
% discipline=CSE
% year=2024
% exam_type=Term Final
% section=A
% ct_number=UNKNOWN
% question_number=2aiii
% topic=Differentiation
% subtopic=Basic Differentiation
% difficulty=Easy
% length=Short
% question_type=New
% frequency=1
% appearances=
% OCR_STATUS=VERIFIED
\section*{Term Final (Sec A) - Q2(a)(iii)}
Question:
Find $\frac{dy}{dx}$ of $x = a \cos^3 t, y = a \sin^3 t$.

Solution:
Step 1: Find the parametric derivatives $\frac{dx}{dt}$ and $\frac{dy}{dt}$.
\[
\frac{dx}{dt} = \frac{d}{dt} (a \cos^3 t) = -3a \cos^2 t \sin t
\]
\[
\frac{dy}{dt} = \frac{d}{dt} (a \sin^3 t) = 3a \sin^2 t \cos t
\]

Step 2: Compute $\frac{dy}{dx}$ using the chain rule.
\[
\frac{dy}{dx} = \frac{\frac{dy}{dt}}{\frac{dx}{dt}} = \frac{3a \sin^2 t \cos t}{-3a \cos^2 t \sin t}
\]
Cancel out common factors:
\[
\frac{dy}{dx} = -\frac{\sin t}{\cos t} = -\tan t
\]

Final Answer:
\[
\boxed{\frac{dy}{dx} = -\tan t}
\]
%%%%%%%%%%%%%%%%%%%%%%%%%%%%%%%%%%%%%%%%%%%%%%%%%%
% QUESTION_END
%%%%%%%%%%%%%%%%%%%%%%%%%%%%%%%%%%%%%%%%%%%%%%%%%%

%%%%%%%%%%%%%%%%%%%%%%%%%%%%%%%%%%%%%%%%%%%%%%%%%%
% QUESTION_START
%%%%%%%%%%%%%%%%%%%%%%%%%%%%%%%%%%%%%%%%%%%%%%%%%%
% id=cse24-final-seca-q2b
% discipline=CSE
% year=2024
% exam_type=Term Final
% section=A
% ct_number=UNKNOWN
% question_number=2b
% topic=Differentiation
% subtopic=Successive Differentiation
% difficulty=Medium
% length=Medium
% question_type=New
% frequency=1
% appearances=
% OCR_STATUS=VERIFIED
\section*{Term Final (Sec A) - Q2(b)}
Question:
Find $n$-th derivatives of $y = e^{2x} \sin(4x+3)$.

Solution:
Step 1: State the general formula for the $n$-th derivative.
For a function of the form $u(x) = e^{ax} \sin(bx+c)$, the $n$-th derivative is:
\[
D^n [e^{ax} \sin(bx+c)] = (a^2+b^2)^{n/2} e^{ax} \sin\left( bx+c + n\phi \right)
\]
where $\tan\phi = \frac{b}{a}$.

Step 2: Identify parameters.
Here, $a = 2$, $b = 4$, and $c = 2$.
Compute the amplitude factor:
\[
\sqrt{a^2 + b^2} = \sqrt{2^2 + 4^2} = \sqrt{4 + 16} = \sqrt{20} = 2\sqrt{5}
\]
Compute the phase angle $\phi$:
\[
\tan\phi = \frac{4}{2} = 2 \implies \phi = \tan^{-1}(2)
\]

Step 3: Formulate the $n$-th derivative.
\[
y_n = (20)^{n/2} e^{2x} \sin\left( 4x+3 + n\tan^{-1}(2) \right)
\]

Final Answer:
\[
\boxed{y_n = 20^{n/2} e^{2x} \sin\left( 4x+3 + n\tan^{-1}(2) \right)}
\]
%%%%%%%%%%%%%%%%%%%%%%%%%%%%%%%%%%%%%%%%%%%%%%%%%%
% QUESTION_END
%%%%%%%%%%%%%%%%%%%%%%%%%%%%%%%%%%%%%%%%%%%%%%%%%%

%%%%%%%%%%%%%%%%%%%%%%%%%%%%%%%%%%%%%%%%%%%%%%%%%%
% QUESTION_START
%%%%%%%%%%%%%%%%%%%%%%%%%%%%%%%%%%%%%%%%%%%%%%%%%%
% id=cse24-final-seca-q2c
% discipline=CSE
% year=2024
% exam_type=Term Final
% section=A
% ct_number=UNKNOWN
% question_number=2c
% topic=Series Expansion
% subtopic=Taylor Series
% difficulty=Medium
% length=Medium
% question_type=New
% frequency=1
% appearances=
% OCR_STATUS=VERIFIED
\section*{Term Final (Sec A) - Q2(c)}
Question:
Expand $\cos x$ in powers of $x - \frac{\pi}{4}$.

Solution:
Step 1: State the Taylor series formula.
The Taylor series expansion of $f(x)$ about $x = a$ is:
\[
f(x) = f(a) + f'(a)(x-a) + \frac{f''(a)}{2!}(x-a)^2 + \frac{f'''(a)}{3!}(x-a)^3 + \dots = \sum_{n=0}^\infty \frac{f^{(n)}(a)}{n!} (x-a)^n
\]
Here, $f(x) = \cos x$ and $a = \frac{\pi}{4}$.

Step 2: Evaluate $f(x)$ and its derivatives at $a = \frac{\pi}{4}$.
\begin{align*}
    f(x) &= \cos x \implies f\left(\frac{\pi}{4}\right) = \frac{1}{\sqrt{2}} \\
    f'(x) &= -\sin x \implies f'\left(\frac{\pi}{4}\right) = -\frac{1}{\sqrt{2}} \\
    f''(x) &= -\cos x \implies f''\left(\frac{\pi}{4}\right) = -\frac{1}{\sqrt{2}} \\
    f'''(x) &= \sin x \implies f'''\left(\frac{\pi}{4}\right) = \frac{1}{\sqrt{2}} \\
    f^{(4)}(x) &= \cos x \implies f^{(4)}\left(\frac{\pi}{4}\right) = \frac{1}{\sqrt{2}}
\end{align*}

Step 3: Construct the expansion.
\[
\cos x = \frac{1}{\sqrt{2}} - \frac{1}{\sqrt{2}}\left(x-\frac{\pi}{4}\right) - \frac{1}{2!\sqrt{2}}\left(x-\frac{\pi}{4}\right)^2 + \frac{1}{3!\sqrt{2}}\left(x-\frac{\pi}{4}\right)^3 + \frac{1}{4!\sqrt{2}}\left(x-\frac{\pi}{4}\right)^4 - \dots
\]
Factoring out $\frac{1}{\sqrt{2}}$:
\[
\cos x = \frac{1}{\sqrt{2}} \left[ 1 - \left(x-\frac{\pi}{4}\right) - \frac{1}{2!}\left(x-\frac{\pi}{4}\right)^2 + \frac{1}{3!}\left(x-\frac{\pi}{4}\right)^3 + \frac{1}{4!}\left(x-\frac{\pi}{4}\right)^4 - \dots \right]
\]

Final Answer:
\[
\boxed{\cos x = \frac{1}{\sqrt{2}} \left[ 1 - \left(x-\frac{\pi}{4}\right) - \frac{1}{2!}\left(x-\frac{\pi}{4}\right)^2 + \frac{1}{3!}\left(x-\frac{\pi}{4}\right)^3 + \frac{1}{4!}\left(x-\frac{\pi}{4}\right)^4 - \dots \right]}
\]
%%%%%%%%%%%%%%%%%%%%%%%%%%%%%%%%%%%%%%%%%%%%%%%%%%
% QUESTION_END
%%%%%%%%%%%%%%%%%%%%%%%%%%%%%%%%%%%%%%%%%%%%%%%%%%

%%%%%%%%%%%%%%%%%%%%%%%%%%%%%%%%%%%%%%%%%%%%%%%%%%
% QUESTION_START
%%%%%%%%%%%%%%%%%%%%%%%%%%%%%%%%%%%%%%%%%%%%%%%%%%
% id=cse24-final-seca-q3a
% discipline=CSE
% year=2024
% exam_type=Term Final
% section=A
% ct_number=UNKNOWN
% question_number=3a
% topic=Differentiation
% subtopic=Leibnitz Rule
% difficulty=Hard
% length=Long
% question_type=New
% frequency=1
% appearances=
% OCR_STATUS=VERIFIED
\section*{Term Final (Sec A) - Q3(a)}
Question:
If $y = (\sinh^{-1} x)^2$, then show that $(1+x^2)y_{n+2} + (2n+1)xy_{n+1} + n^2 y_n = 0$, also find $(y_n)_0$.

Solution:
Step 1: Compute first and second derivatives.
Given $y = (\sinh^{-1} x)^2$. Differentiating once with respect to $x$:
\[
y_1 = 2(\sinh^{-1} x) \cdot \frac{1}{\sqrt{1+x^2}}
\]
Multiply both sides by $\sqrt{1+x^2}$ and square:
\[
(1+x^2)y_1^2 = 4(\sinh^{-1} x)^2 \implies (1+x^2)y_1^2 = 4y
\]
Differentiate both sides again:
\[
2x y_1^2 + (1+x^2)\cdot 2y_1y_2 = 4y_1
\]
Divide by $2y_1$ (since $y_1 \ne 0$):
\[
(1+x^2)y_2 + x y_1 = 2 \quad \text{--- (Equation 1)}
\]

Step 2: Differentiate Equation 1 $n$ times using Leibnitz's theorem.
\[
D^n\left[ (1+x^2)y_2 \right] + D^n[x y_1] = 0
\]
- For $D^n\left[(1+x^2)y_2\right]$:
  \[
  (1+x^2)y_{n+2} + n(2x)y_{n+1} + \frac{n(n-1)}{2}(2)y_n = (1+x^2)y_{n+2} + 2n x y_{n+1} + n(n-1)y_n
  \]
- For $D^n[x y_1]$:
  \[
  x y_{n+1} + n(1)y_n = x y_{n+1} + n y_n
  \]
Summing both components:
\[
(1+x^2)y_{n+2} + (2n x + x)y_{n+1} + [n(n-1) + n]y_n = 0
\]
\[
(1+x^2)y_{n+2} + (2n+1)xy_{n+1} + n^2 y_n = 0
\]
This proves the first required equation.

Step 3: Evaluate derivatives at $x = 0$.
- At $x = 0$:
  \[
  y(0) = (\sinh^{-1} 0)^2 = 0
  \]
  \[
  y_1(0) = 0
  \]
  From Equation 1:
  \[
  (1+0)y_2(0) + 0 \cdot y_1(0) = 2 \implies y_2(0) = 2
  \]
  From the general recurrence relation at $x = 0$:
  \[
  (y_{n+2})_0 = -n^2 (y_n)_0
  \]

Step 4: Resolve the sequence $(y_n)_0$ based on odd/even values of $n$.
- \textbf{If $n$ is odd}:
  Since $y_1(0) = 0$, all odd-ordered derivatives evaluated at $0$ are zero:
  \[
  (y_n)_0 = 0 \quad (\text{for odd } n)
  \]
- \textbf{If $n$ is even} (let $n = 2k$):
  \[
  (y_{2k})_0 = -(2k-2)^2 (y_{2k-2})_0 = (-1)^2 (2k-2)^2 (2k-4)^2 (y_{2k-4})_0 = \dots
  \]
  \[
  (y_{2k})_0 = (-1)^{k-1} [ (2k-2)(2k-4)\dots 2 ]^2 (y_2)_0 = (-1)^{k-1} \left[ 2^{k-1}(k-1)! \right]^2 \cdot 2 = (-1)^{k-1} 2^{2k-1} \left[ (k-1)! \right]^2
  \]

Final Answer:
\[
\boxed{
(y_n)_0 = \begin{cases} 
0, & \text{if } n \text{ is odd} \\ 
(-1)^{\frac{n}{2}-1} 2^{n-1} \left[ \left(\frac{n}{2}-1\right)! \right]^2, & \text{if } n \text{ is even (where } n \ge 2)
\end{cases}
}
\]
%%%%%%%%%%%%%%%%%%%%%%%%%%%%%%%%%%%%%%%%%%%%%%%%%%
% QUESTION_END
%%%%%%%%%%%%%%%%%%%%%%%%%%%%%%%%%%%%%%%%%%%%%%%%%%

%%%%%%%%%%%%%%%%%%%%%%%%%%%%%%%%%%%%%%%%%%%%%%%%%%
% QUESTION_START
%%%%%%%%%%%%%%%%%%%%%%%%%%%%%%%%%%%%%%%%%%%%%%%%%%
% id=cse24-final-seca-q3b
% discipline=CSE
% year=2024
% exam_type=Term Final
% section=A
% ct_number=UNKNOWN
% question_number=3b
% topic=Applications of Derivatives
% subtopic=Maxima \& Minima
% difficulty=Hard
% length=Long
% question_type=New
% frequency=1
% appearances=
% OCR_STATUS=VERIFIED
\section*{Term Final (Sec A) - Q3(b)}
Question:
Distinguish between critical points and points of inflection of a function. Find the extremum values of $f(x) = e^x + 2\cos x + e^{-x}$.

Solution:
Step 1: Distinguish between critical points and points of inflection.
\begin{itemize}
    \item \textbf{Critical Points}: Points $c$ in the domain of $f(x)$ where $f'(c) = 0$ or $f'(c)$ is undefined. These represent potential locations for local extrema (maxima/minima).
    \item \textbf{Points of Inflection}: Points on a curve where the concavity changes. At these points, $f''(x) = 0$ (or is undefined) and $f''(x)$ changes its sign.
\end{itemize}

Step 2: Find critical points of $f(x) = e^x + 2\cos x + e^{-x}$.
We can rewrite $f(x) = 2\cosh x + 2\cos x$.
Differentiate with respect to $x$:
\[
f'(x) = e^x - 2\sin x - e^{-x} = 2\sinh x - 2\sin x
\]
Set $f'(x) = 0 \implies \sinh x = \sin x$.
Since $\sinh x > x > \sin x$ for $x > 0$ and $\sinh x < x < \sin x$ for $x < 0$, the only real solution to $\sinh x = \sin x$ is:
\[
x = 0
\]

Step 3: Analyze the nature of the critical point $x = 0$ using higher-order derivatives.
Compute successive derivatives at $x = 0$:
\begin{align*}
    f''(x) &= e^x - 2\cos x + e^{-x} \implies f''(0) = 1 - 2 + 1 = 0 \\
    f'''(x) &= e^x + 2\sin x - e^{-x} = 2\sinh x + 2\sin x \implies f'''(0) = 0 \\
    f^{(4)}(x) &= e^x + 2\cos x + e^{-x} \implies f^{(4)}(0) = 1 + 2 + 1 = 4 > 0
\end{align*}
Since the first non-zero derivative at the stationary point is of even order (4th) and is positive ($4 > 0$), $x = 0$ is a point of local minimum.

Step 4: Compute the extremum (minimum) value.
\[
f(0) = e^0 + 2\cos 0 + e^{-0} = 1 + 2 + 1 = 4
\]

Final Answer:
\[
\boxed{\text{The function has a minimum value of } 4 \text{ at } x = 0.}
\]
%%%%%%%%%%%%%%%%%%%%%%%%%%%%%%%%%%%%%%%%%%%%%%%%%%
% QUESTION_END
%%%%%%%%%%%%%%%%%%%%%%%%%%%%%%%%%%%%%%%%%%%%%%%%%%

%%%%%%%%%%%%%%%%%%%%%%%%%%%%%%%%%%%%%%%%%%%%%%%%%%
% QUESTION_START
%%%%%%%%%%%%%%%%%%%%%%%%%%%%%%%%%%%%%%%%%%%%%%%%%%
% id=cse24-final-seca-q4a
% discipline=CSE
% year=2024
% exam_type=Term Final
% section=A
% ct_number=UNKNOWN
% question_number=4a
% topic=Applications of Derivatives
% subtopic=Tangent \& Normal
% difficulty=Hard
% length=Long
% question_type=New
% frequency=1
% appearances=
% OCR_STATUS=VERIFIED
\section*{Term Final (Sec A) - Q4(a)}
Question:
Define tangent, subtangent, normal and subnormal of a curve $f(x)=0$ at $x=x_0$. Find the equation of tangent and normal line to the curve $x=e^{-t}\cos t, y=e^t\sin t$ at $t=\pi$.

Solution:
Step 1: Geometrical Definitions.
Let $P(x_0, y_0)$ be a point on a curve. Draw a tangent and normal at $P$ that intersect the $x$-axis at $T$ and $N$ respectively. Let $G(x_0, 0)$ be the foot of the perpendicular from $P$ to the $x$-axis.
\begin{itemize}
    \item \textbf{Tangent Length}: The length of the segment $PT$ on the tangent line.
    \item \textbf{Subtangent}: The length of the projection segment $TG$ on the $x$-axis.
    \item \textbf{Normal Length}: The length of the segment $PN$ on the normal line.
    \item \textbf{Subnormal}: The length of the projection segment $GN$ on the $x$-axis.
\end{itemize}

Step 2: Find the coordinates of the point at $t = \pi$.
\[
x_0 = e^{-\pi} \cos\pi = -e^{-\pi}
\]
\[
y_0 = e^\pi \sin\pi = 0
\]
So, the point of contact is $P\left(-e^{-\pi}, 0\right)$.

Step 3: Find the derivative $\frac{dy}{dx}$ at $t = \pi$.
Differentiate the parametric equations:
\[
\frac{dx}{dt} = -e^{-t}\cos t - e^{-t}\sin t = -e^{-t}(\cos t + \sin t)
\]
\[
\frac{dy}{dt} = e^t\sin t + e^t\cos t = e^t(\sin t + \cos t)
\]
Using the chain rule:
\[
\frac{dy}{dx} = \frac{\frac{dy}{dt}}{\frac{dx}{dt}} = \frac{e^t(\sin t + \cos t)}{-e^{-t}(\cos t + \sin t)} = -e^{2t}
\]
At $t = \pi$:
\[
m = \left.\frac{dy}{dx}\right|_{t=\pi} = -e^{2\pi}
\]

Step 4: Find the equations of the tangent and normal.
\begin{itemize}
    \item \textbf{Tangent Line}:
    \[
    y - y_0 = m(x - x_0) \implies y - 0 = -e^{2\pi}\left(x - \left(-e^{-\pi}\right)\right)
    \]
    \[
    y = -e^{2\pi} x - e^{\pi} \implies e^{2\pi}x + y + e^\pi = 0
    \]
    \item \textbf{Normal Line}:
    The slope of the normal is $-\frac{1}{m} = e^{-2\pi}$:
    \[
    y - y_0 = -\frac{1}{m}(x - x_0) \implies y - 0 = e^{-2\pi}\left(x - \left(-e^{-\pi}\right)\right)
    \]
    \[
    y = e^{-2\pi} x + e^{-3\pi} \implies e^{-2\pi}x - y + e^{-3\pi} = 0 \quad (\text{or } e^\pi x - e^{3\pi} y + 1 = 0)
    \]
\end{itemize}

Final Answer:
\[
\boxed{
\begin{aligned}
\text{Tangent: } & e^{2\pi}x + y + e^\pi = 0 \\
\text{Normal: } & e^\pi x - e^{3\pi}y + 1 = 0
\end{aligned}
}
\]
%%%%%%%%%%%%%%%%%%%%%%%%%%%%%%%%%%%%%%%%%%%%%%%%%%
% QUESTION_END
%%%%%%%%%%%%%%%%%%%%%%%%%%%%%%%%%%%%%%%%%%%%%%%%%%

%%%%%%%%%%%%%%%%%%%%%%%%%%%%%%%%%%%%%%%%%%%%%%%%%%
% QUESTION_START
%%%%%%%%%%%%%%%%%%%%%%%%%%%%%%%%%%%%%%%%%%%%%%%%%%
% id=cse24-final-seca-q4b
% discipline=CSE
% year=2024
% exam_type=Term Final
% section=A
% ct_number=UNKNOWN
% question_number=4b
% topic=Limits
% subtopic=General
% difficulty=Medium
% length=Long
% question_type=New
% frequency=1
% appearances=
% OCR_STATUS=VERIFIED
\section*{Term Final (Sec A) - Q4(b)}
Question:
Describe the indeterminate forms of limits. Evaluate the following using L'Hospital rule:
(i) $\lim_{x \to 0} \frac{3 \tan x - 3x - x^3}{x^5}$ \quad
(ii) $\lim_{x \to 0} (\sin x)^x$

Solution:
Step 1: Describe the Indeterminate Forms.
In calculus, an indeterminate form is an algebraic expression obtained when evaluating limits whose limit cannot be determined solely from the limits of the individual parts. The seven classic indeterminate forms are:
\[
\frac{0}{0}, \quad \frac{\infty}{\infty}, \quad 0 \cdot \infty, \quad \infty - \infty, \quad 0^0, \quad 1^\infty, \quad \infty^0
\]

Step 2: Evaluate (i) $\lim_{x \to 0} \frac{3 \tan x - 3x - x^3}{x^5}$.
The limit initially presents the indeterminate form $\frac{0}{0}$. Differentiating repeatedly using L'Hospital's rule:
\[
L_1 = \lim_{x \to 0} \frac{3\sec^2 x - 3 - 3x^2}{5x^4} \quad \left[\frac{0}{0}\right]
\]
Differentiate again:
\[
L_1 = \lim_{x \to 0} \frac{6\sec^2 x \tan x - 6x}{20x^3} \quad \left[\frac{0}{0}\right]
\]
Differentiate again:
\[
L_1 = \lim_{x \to 0} \frac{6\left(\sec^4 x + 2\sec^2 x \tan^2 x\right) - 6}{60x^2} \quad \left[\frac{0}{0}\right]
\]
Substitute $\sec^2 x = 1 + \tan^2 x$ to simplify:
\[
\sec^4 x + 2\sec^2 x \tan^2 x - 1 = (1 + \tan^2 x)^2 + 2(1 + \tan^2 x)\tan^2 x - 1 = 4\tan^2 x + 3\tan^4 x
\]
Thus:
\[
L_1 = \lim_{x \to 0} \frac{6\left(4\tan^2 x + 3\tan^4 x\right)}{60x^2} = \lim_{x \to 0} \frac{24\tan^2 x + 18\tan^4 x}{60x^2}
\]
Since $\lim_{x \to 0} \frac{\tan x}{x} = 1$:
\[
L_1 = \lim_{x \to 0} \left[ \frac{24}{60} \left(\frac{\tan x}{x}\right)^2 + \frac{18}{60} \frac{\tan^4 x}{x^2} \right] = \frac{24}{60} + 0 = \frac{2}{5}
\]

Step 3: Evaluate (ii) $\lim_{x \to 0} (\sin x)^x$.
This presents the indeterminate form $0^0$. Let $L_2 = \lim_{x \to 0} (\sin x)^x$, then:
\[
\ln L_2 = \lim_{x \to 0} x \ln(\sin x) = \lim_{x \to 0} \frac{\ln(\sin x)}{\frac{1}{x}} \quad \left[\frac{-\infty}{\infty}\right]
\]
Applying L'Hospital's rule:
\[
\ln L_2 = \lim_{x \to 0} \frac{\cot x}{-\frac{1}{x^2}} = \lim_{x \to 0} \frac{-x^2}{\tan x} = \lim_{x \to 0} (-x) \cdot \left(\frac{x}{\tan x}\right) = 0 \cdot 1 = 0
\]
Thus, $L_2 = e^0 = 1$.

Final Answer:
\[
\boxed{\text{(i) } \frac{2}{5}, \quad \text{(ii) } 1}
\]
%%%%%%%%%%%%%%%%%%%%%%%%%%%%%%%%%%%%%%%%%%%%%%%%%%
% QUESTION_END
%%%%%%%%%%%%%%%%%%%%%%%%%%%%%%%%%%%%%%%%%%%%%%%%%%

%%%%%%%%%%%%%%%%%%%%%%%%%%%%%%%%%%%%%%%%%%%%%%%%%%
% QUESTION_START
%%%%%%%%%%%%%%%%%%%%%%%%%%%%%%%%%%%%%%%%%%%%%%%%%%
% id=cse24-final-seca-q4c
% discipline=CSE
% year=2024
% exam_type=Term Final
% section=A
% ct_number=UNKNOWN
% question_number=4c
% topic=Partial Derivatives
% subtopic=Total Derivatives
% difficulty=Medium
% length=Medium
% question_type=New
% frequency=1
% appearances=
% OCR_STATUS=VERIFIED
\section*{Term Final (Sec A) - Q4(c)}
Question:
If $u = \ln(x^2+y^2+z^2)$ then construct the relation $x\frac{\partial u}{\partial x} + y\frac{\partial u}{\partial y} + z\frac{\partial u}{\partial z} = 2$.

Solution:
Step 1: Find the partial derivatives with respect to $x$, $y$, and $z$.
Using the chain rule:
\[
\frac{\partial u}{\partial x} = \frac{1}{x^2+y^2+z^2} \cdot \frac{\partial}{\partial x}(x^2+y^2+z^2) = \frac{2x}{x^2+y^2+z^2}
\]
Similarly, by symmetry:
\[
\frac{\partial u}{\partial y} = \frac{2y}{x^2+y^2+z^2}
\]
\[
\frac{\partial u}{\partial z} = \frac{2z}{x^2+y^2+z^2}
\]

Step 2: Substitute these partial derivatives into the target expression.
\[
x\frac{\partial u}{\partial x} + y\frac{\partial u}{\partial y} + z\frac{\partial u}{\partial z} = x\left(\frac{2x}{x^2+y^2+z^2}\right) + y\left(\frac{2y}{x^2+y^2+z^2}\right) + z\left(\frac{2z}{x^2+y^2+z^2}\right)
\]
\[
= \frac{2x^2 + 2y^2 + 2z^2}{x^2+y^2+z^2}
\]
Factor out $2$ from the numerator:
\[
= \frac{2(x^2+y^2+z^2)}{x^2+y^2+z^2} = 2
\]
This successfully establishes and constructs the required relation.

Final Answer:
\[
\boxed{x\frac{\partial u}{\partial x} + y\frac{\partial u}{\partial y} + z\frac{\partial u}{\partial z} = 2}
\]
%%%%%%%%%%%%%%%%%%%%%%%%%%%%%%%%%%%%%%%%%%%%%%%%%%
% QUESTION_END
%%%%%%%%%%%%%%%%%%%%%%%%%%%%%%%%%%%%%%%%%%%%%%%%%%

\end{document}